\worldchapter{Magie und Glaube}{magic}
\label{ch:magie}

\begin{multicols*}{2}
\section{Zwei ubernaturliche Motoren}
Das Spiel trennt bewusst zwischen:
\begin{itemize}[leftmargin=*]
\item \textbf{Zauberei}: wiederholbare, technische Zauberei
\item \textbf{Belief}: Bitte, Ritus und Antwort hoherer Machte
\end{itemize}

Beide sind gefahrlich und teuer, funktionieren aber nicht gleich.

\section{Zauberei}
\subsection{Auslöser}
Zauberei wird genutzt, wenn ein Charakter okkulte Kraft direkt formt.

\subsection{Zugriff}
\begin{itemize}[leftmargin=*]
\item Zauberei >= 1: eine Schule wahlen
\item Zauberei >= 2: zwei Schulen wahlen
\item Startzauber: 2 + Zauberei Rang, maximal 5
\item bekannte Zauber mussen aus den gewahlten Schulen stammen
\end{itemize}

Aktive Schulen sind:
\begin{itemize}[leftmargin=*]
\item Wizardry
\item White Magic
\item Black Magic
\item Elemental Magic
\item Sanguine Magic
\item Astral Magic
\item Lizard Folk Magic
\end{itemize}

\subsection{Wirkvorgang}
\begin{enumerate}[leftmargin=*]
\item einen bekannten Zauber wahlen
\item Wirkmodus wahlen
\item Stufe wahlen
\item Mind oder Will + Zauberei Probe
\item Arcana-Pool-Kosten immer bezahlen, auch bei Fehlschlag
\end{enumerate}

\subsection{Wirkmodi}
\begin{itemize}[leftmargin=*]
\item \textbf{Fast Cast}: 1 Action
\item \textbf{Ritual Cast}: 2+ Actions oder ein Downtime-Segment, wenn Vorbereitung, Kreise, Hilfen oder Material notig sind
\end{itemize}

\subsection{Wirkstufen}
\begin{description}[style=nextline,leftmargin=0pt]
\item[\textbf{Minor}] +10 Zielwert, kein Zuschlag
\item[\textbf{Severe}] -10 Zielwert, +1 Arcana
\item[\textbf{Catastrophic}] -30 Zielwert, +2 Arcana
\end{description}

\subsection{Arcana-Kosten}
Basiskosten nach spell\_points:
\begin{itemize}[leftmargin=*]
\item 1--3 \textrightarrow{} 1 Arcana
\item 4--7 \textrightarrow{} 2 Arcana
\item 8--11 \textrightarrow{} 3 Arcana
\item 12+ \textrightarrow{} 4 Arcana
\end{itemize}

Gesamtkosten sind Basis plus Stufenzuschlag, Minimum 1.

\subsection{Ergebnisse}
\begin{itemize}[leftmargin=*]
\item Kritischer, starker oder sauberer Erfolg: Effekt tritt wie erklart ein
\item Knapp geschafft: Effekt tritt ein, aber mit einem kleinen instabilen Detail
\item jeder Misserfolg: kein beabsichtigter Effekt, stattdessen Backlash
\end{itemize}

\subsection{Zauberverstärkung}
Bei erfolgreicher Zauberei gibt \textbf{Magic Level} genau \textbf{einen} Boost-Schritt auf:
\begin{itemize}[leftmargin=*]
\item Reichweite
\item Dauer
\item Zielzahl
\item Schaden- oder Clear-Wert
\end{itemize}

\subsection{Backlash}
\begin{itemize}[leftmargin=*]
\item \textbf{Minor}: +1 Burden oder \textbf{Shaken}
\item \textbf{Serious}: +1 Wound, groBer Kollateralschaden oder feindliche Aufmerksamkeit
\item \textbf{Critical}: +2 Burden, schwerer Kollateralschaden, Riss oder bleibender Fluchmal-Effekt
\end{itemize}

Bei einem misslungenen 96--00 kommt zusatzlich ein thematischer Nebeneffekt hinzu.

\subsection{Chaotische Zauberei}
Chaotische Zauberei gilt, wenn:
\begin{itemize}[leftmargin=*]
\item Magic Level >= 5, oder
\item der zaubernde Charakter 2+ Verderbnis hat
\end{itemize}

Nach jedem Wirken 1d6:
\begin{enumerate}[leftmargin=*]
\item Echo: Effekt trifft ein nachstes gultiges Nebenziel mit Limited Effect
\item Drain: 1 zusatzliches Arcana verlieren
\item Fracture: nahe Zerbrechlichkeit oder Ward bricht
\item Beacon: okkulte Aufmerksamkeit richtet sich auf die Zaubernden
\item Reversal Pulse: ein kleines beabsichtigtes Detail kehrt sich um
\item Stable Surge: kein Zusatzeffekt
\end{enumerate}

\section{Belief}
\subsection{Auslöser}
Belief wird genutzt, wenn ein Charakter eine Macht, Heilige, Ahnen oder namenlose Kräfte um Eingriff bittet.

\subsection{Petition-Ablauf}
\begin{enumerate}[leftmargin=*]
\item Bitte, Zeugin und Opfer benennen
\item \textbf{Gunst} ausgeben
\item optional +1 Burden fur eine verzweifelte Bitte nehmen
\item Petition werfen:
\[
1d10 + Anrufungswert + Gunst + Modifikatoren
\]
\end{enumerate}

\[
\text{Anrufungswert} = \text{Faith Level} + \text{Faith Rang}
\]

Maximale Gunst pro Anrufung:
\[
\text{Maximale Gunst pro Anrufung} = 1 + \lfloor \text{Will} / 2 \rfloor
\]

\subsection{Zeremonielle Modifikatoren}
\begin{itemize}[leftmargin=*]
\item +1 passendes Opfer, Gelubde oder Preis
\item +1 geweihter Ort, Relikt oder anerkannte Zeugin
\item +1 Chor, Helfende oder Gemeinde
\item -1 hastig, improvisiert oder unterbrochen
\item -2 Bitte widerspricht klar Lehre, Eid oder Natur der angerufenen Macht
\item Gesamtkappung: hochstens +2, niedrigstens -2
\end{itemize}

\subsection{Ergebnisbander}
\begin{description}[style=nextline,leftmargin=0pt]
\item[\textbf{1--7 Silence}] keine direkte Intervention, nur Zeichen oder Schweigen
\item[\textbf{8--11 Sign}] symbolische Antwort, Warnung, Hinweis oder kurzer Schutz
\item[\textbf{12--15 Intercession}] spurbare Wunderwirkung im Rahmen der Szene
\item[\textbf{16+ Theophany}] groBer Eingriff, fremd und verstorend
\end{description}

Gunst wird immer verbraucht. Bei \textbf{Silence} unter Tabubruch gibt es +1 Burden. Bei \textbf{Theophany} wahlt der Guide einen heiligen Preis: Gelubde, Stigma oder Omen-Schuld.

\section{Petitionsformen}
\subsection{Shock Prayer}
\begin{itemize}[leftmargin=*]
\item Ausloser: unmittelbare Gefahr
\item Dauer: 1 Action
\item Kosten: 1--2 Gunst
\item Wirkung: unmittelbare Standfestigkeit, Flucht, Warnung, Schutz oder Wucht
\end{itemize}
Bei \textbf{Sign}: +Faith Level auf einen sofortigen Test oder eine eingehende Consequence um eine Stufe senken.\\
Bei \textbf{Intercession}: zwei solche Boons.\\
Bei \textbf{Theophany}: die Szene kippt deutlich zugunsten der Gruppe, aber der Preis folgt.

\subsection{Blessing Rite}
\begin{itemize}[leftmargin=*]
\item Dauer: 10 Minuten
\item Kosten: 2--4 Gunst
\item Ziel: Person, Objekt, Schwelle, Grab, Wasserstelle oder kleine gottliche Verunreinigung
\end{itemize}
Sign wirkt bis zum Szenenende, Intercession bis zur nachsten Full Rest, Theophany kann hartnackige Fluche brechen oder heiligen Boden schaffen.

\subsection{Lesser Request}
\begin{itemize}[leftmargin=*]
\item Dauer: 10--15 Minuten
\item Kosten: 2--6 Gunst
\item Ziel: eine begrenzte Wunderbitte wie enthullen, schutzen, fuhren, reinigen, beruhigen oder plagen
\end{itemize}
Bei \textbf{Sign} gibt es ein klares Omen, eine Warnung, einen Weg oder kurzen Schutz.\\
Bei \textbf{Intercession} geschieht ein konkretes Wunder auf SzenenmaBstab.\\
Bei \textbf{Theophany} kippt die ganze Szene zugunsten der Bittenden, aber in verstorender oder unheimlicher Form.

\subsection{Great Invocation}
\begin{itemize}[leftmargin=*]
\item Dauer: 30+ Minuten
\item Kosten: 4--8 Gunst
\item Voraussetzung: mindestens eine Zeugin und ein bedeutendes Opfer, Relikt oder Eid
\end{itemize}
Dies ist die Form fur offentliche Rettung, Urteil, Massenprotektion oder heiliges Strafgericht.
Bei \textbf{Sign} kommen Furchtzeichen, Gewissheit oder ein zwingender Auftrag uber die Szene.\\
Bei \textbf{Intercession} greift ein groBes Wunder uber Menge, Stadtteilrand, Schlachtfeldkante oder entweihten Ort.\\
Bei \textbf{Theophany} manifestiert sich eine szenenbestimmende Rettung, Verwandlung oder Verdammnis und fordert trotzdem einen heiligen Preis.

\section{Verbotene Praxis und Relikte}
\begin{itemize}[leftmargin=*]
\item \textbf{Relikt}: speichert 1--3 Arcana oder Gunst und ubertragt sie durch Beruhrung oder Ritus
\item \textbf{Verbotene Methode}: senkt Zauberkosten um 1, mindestens auf 1, verursacht aber 1 Corruption
\end{itemize}

\section{Vollstandiger Zauberkatalog}
Der komplette aktive Katalog aus 178 Zaubern steht in \cref{app:spells}. Fur die Laufzeit am Tisch sind dort vor allem id, name, school, spell\_points und description verbindlich.
\end{multicols*}
